\documentclass[../main.tex]{subfiles}

\begin{document}

\section{Minimos Quadrados Polinomial}

\subsection{Caso discreto}

No caso discreto, a entrada e um conjunto de pontos \((x_i, y_i)\). Para ajustar
um polinomio de grau \(m\),
\[
    p(x) = c_0 + c_1x + c_2x^2 + \cdots + c_mx^m,
\]
o codigo monta as equacoes normais usando somatorios. A matriz do sistema recebe
os termos \(\sum x_i^{j+k}\), enquanto o vetor independente recebe os termos
\(\sum y_i x_i^k\).

Depois de montar o sistema, os coeficientes sao encontrados por eliminacao de
Gauss implementada no proprio projeto.

\subsection{Caso continuo}

No caso continuo, o objetivo e aproximar uma funcao \(f(x)\) em um intervalo
\([a,b]\). As equacoes normais usam integrais:
\[
    \int_a^b p(x)x^k dx = \int_a^b f(x)x^k dx.
\]

As integrais de potencias de \(x\), que formam a matriz, sao calculadas pela
formula exata. Ja as integrais que envolvem \(f(x)\) sao aproximadas pela regra
composta de Simpson, tambem implementada diretamente no codigo.

\subsection{Arquivos de entrada e saida}

O caso discreto le \path{input/mmq_discreto.txt}, com um ponto por linha. O caso
continuo le o arquivo \path{input/mmq_continuo.txt}. Nele sao informados os
parametros \texttt{funcao}, \texttt{a}, \texttt{b}, \texttt{grau} e
\texttt{subintervalos}.

Ambos os metodos escrevem CSV em UTF-8 na pasta \texttt{output}, contendo o polinomio
encontrado, os coeficientes e uma tabela de comparacao entre os valores reais e
os valores aproximados.

\subsection{Problemas teste}

\subsubsection{Exercicio 8.1: regressao quadratica da taxa de acidentes}

A segunda parte do exercicio 8.1 pede uma regressao quadratica para os acidentes
por 10000 veiculos. Novamente foi usado \(x = ano - 1980\). O polinomio obtido
foi:
\[
    p(x) = 1.6985 - 0.0368563x + 0.00115007x^2.
\]
Para o ano 2000, isto e, \(x=20\), o valor encontrado foi:
\[
    p(20) = 1.42140.
\]
Esse resultado representa a previsao de aproximadamente \(1.4214\) acidentes por
10000 veiculos. O CSV gerado esta em
\texttt{output/exercicios/ex8\_1.csv}.

\subsubsection{Exercicio 8.5: coeficiente da placa de orificio}

O exercicio 8.5 fornece valores experimentais do coeficiente \(C\) em funcao de
\(x=A/A_1\), e pede o ajuste por
\[
    C(x) = a_0 + a_1x + a_2x^2.
\]
Com os dados de \texttt{input/exercicios/ex8\_5\_orificio.txt}, o programa
obteve:
\[
    C(x) = 0.650167 - 0.231742x + 0.564394x^2.
\]

Alguns valores calculados foram:
\[
\begin{array}{c|c|c}
x & C\text{ da tabela} & C\text{ ajustado} \\
\hline
0.10 & 0.62 & 0.632636 \\
0.50 & 0.68 & 0.675394 \\
0.80 & 0.81 & 0.825985 \\
1.00 & 1.00 & 0.982818
\end{array}
\]
Os residuos ficaram pequenos para a maioria dos pontos, indicando um ajuste
coerente com a tendencia crescente dos dados. A tabela completa esta em
\path{output/exercicios/ex8_5.csv}.

\subsubsection{Exercicio 8.11: PNB por modelo de potencia}

O exercicio 8.11 pede o ajuste dos PNBs por uma funcao do tipo \(a x^b\). Para
usar minimos quadrados, apliquei logaritmo nos dois lados:
\[
    \ln y = \ln a + b\ln x.
\]
Como o modelo de potencia exige \(x>0\), usei \(x=ano-1979\).

Para o PNB em dolares correntes, o programa obteve:
\[
    y = 225.107890x^{0.425689},
\]
com previsao para 2000 igual a \(822.7053\) milhoes. Para o PNB em dolares
constantes:
\[
    y = 329.451061x^{0.236483},
\]
com previsao para 2000 igual a \(676.8197\) milhoes.

A razao entre os valores previstos foi \(1.21555\). Isso indica que, no ano
2000, o PNB em dolares correntes ficaria cerca de \(21.55\%\) acima do valor em
dolares constantes segundo o modelo ajustado. Os resultados foram salvos em
\path{output/exercicios/ex8_11.csv}.

\end{document}
